\documentclass[10pt,a4paper]{amsart}
\usepackage{amsmath,amssymb,amsthm,graphicx}
\usepackage{verbatim}
\newenvironment{reference}{\comment}{\endcomment}
%\newenvironment{reference}{}{}
\newenvironment{slogan}{\comment}{\endcomment}
\newenvironment{history}{\comment}{\endcomment}

% For commutative diagrams we use Xy-pic
\usepackage[all]{xy}

% We use 2cell for 2-commutative diagrams.
\xyoption{2cell}
% Theorem environments.
%
\theoremstyle{plain}
\newtheorem{theorem}[subsection]{Theorem}
\newtheorem{proposition}[subsection]{Proposition}
\newtheorem{lemma}[subsection]{Lemma}

\theoremstyle{definition}
\newtheorem{definition}[subsection]{Definition}
\newtheorem{example}[subsection]{Example}
\newtheorem{exercise}[subsection]{Exercise}
\newtheorem{situation}[subsection]{Situation}

\theoremstyle{remark}
\newtheorem{remark}[subsection]{Remark}
\newtheorem{remarks}[subsection]{Remarks}

\numberwithin{equation}{subsection}

% Macros
%
\def\lim{\mathop{\mathrm{lim}}\nolimits}
\def\colim{\mathop{\mathrm{colim}}\nolimits}
\def\Spec{\mathop{\mathrm{Spec}}}
\def\Hom{\mathop{\mathrm{Hom}}\nolimits}
\def\Ext{\mathop{\mathrm{Ext}}\nolimits}
\def\SheafHom{\mathop{\mathcal{H}\!\mathit{om}}\nolimits}
\def\SheafExt{\mathop{\mathcal{E}\!\mathit{xt}}\nolimits}
\def\Sch{\mathit{Sch}}
\def\Mor{\mathop{\mathrm{Mor}}\nolimits}
\def\Ob{\mathop{\mathrm{Ob}}\nolimits}
\def\Sh{\mathop{\mathit{Sh}}\nolimits}
\def\NL{\mathop{N\!L}\nolimits}
\def\CH{\mathop{\mathrm{CH}}\nolimits}
\def\proetale{{pro\text{-}\acute{e}tale}}
\def\etale{{\acute{e}tale}}
\def\QCoh{\mathit{QCoh}}
\def\Ker{\mathop{\mathrm{Ker}}}
\def\Im{\mathop{\mathrm{Im}}}
\def\Coker{\mathop{\mathrm{Coker}}}
\def\Coim{\mathop{\mathrm{Coim}}}

% Boxtimes
%
\DeclareMathSymbol{\boxtimes}{\mathbin}{AMSa}{"02}

%
% Macros for moduli stacks/spaces
%
\def\QCohstack{\mathcal{QC}\!\mathit{oh}}
\def\Cohstack{\mathcal{C}\!\mathit{oh}}
\def\Spacesstack{\mathcal{S}\!\mathit{paces}}
\def\Quotfunctor{\mathrm{Quot}}
\def\Hilbfunctor{\mathrm{Hilb}}
\def\Curvesstack{\mathcal{C}\!\mathit{urves}}
\def\Polarizedstack{\mathcal{P}\!\mathit{olarized}}
\def\Complexesstack{\mathcal{C}\!\mathit{omplexes}}
% \Pic is the operator that assigns to X its picard group, usage \Pic(X)
% \Picardstack_{X/B} denotes the Picard stack of X over B
% \Picardfunctor_{X/B} denotes the Picard functor of X over B
\def\Pic{\mathop{\mathrm{Pic}}\nolimits}
\def\Picardstack{\mathcal{P}\!\mathit{ic}}
\def\Picardfunctor{\mathrm{Pic}}
\def\Deformationcategory{\mathcal{D}\!\mathit{ef}}
\usepackage{cleveref}
\crefname{originallemma}{original lemma}{original lemmas}
\Crefname{originallemma}{Original lemma}{Original lemmas}
\makeatletter
\newlabel{more-algebra-lemma-henselization}{{0A02}{source}{}{sourcepointer-1}{}}
\newlabel{more-algebra-lemma-henselization@cref}{{[originallemma][0][]0A02}{[1][0][]source}{}{}{}}
\makeatother
\crefname{originallemma}{original lemma}{original lemmas}
\Crefname{originallemma}{Original lemma}{Original lemmas}
\makeatletter
\newlabel{more-algebra-lemma-henselization-Noetherian-pair}{{0AGV}{source}{}{sourcepointer-2}{}}
\newlabel{more-algebra-lemma-henselization-Noetherian-pair@cref}{{[originallemma][0][]0AGV}{[1][0][]source}{}{}{}}
\makeatother
\title{Completed quotients of henselizations come from finite type algebras}
\author{The Stacks Project Authors}
\date{stacks-project:restricted.tex:local-snapshot-20261012; GNU FDL 1.2 or later}
\begin{document}
\maketitle
\noindent\textit{Editorial note.} All delivered mathematical text is the author's own, copied from the pinned Stacks Project chapter file restricted.tex; no statement or proof step was rewritten, reordered or removed.\par\medskip
\noindent\textit{Editorial note.} The counted proof is the author's complete argument: The easy half of the algebraization toolkit: it produces the finite type model whose formal completion is a prescribed completed quotient of a henselization.\par\medskip
\noindent\textit{Editorial note.} The label lemma-algebraize-easy names the counted statement, and the proof environment immediately following it in the source is delivered verbatim as its proof.\par\medskip
\noindent\textit{Editorial note.} The counted statement is carried by the author's own lemma environment in the chapter source.\par\medskip
\noindent\textit{Allowed standard prerequisites (editorial).}\begin{itemize}
\item More on Algebra, Lemma more-algebra-lemma-henselization-Noetherian-pair: the henselization of a finite type algebra over a Noetherian ring with respect to an ideal is Noetherian
\item More on Algebra, Lemma more-algebra-lemma-henselization (construction): the henselization is a filtered colimit of etale C-algebras C-prime with C/IC isomorphic to C-prime/IC-prime
\item Standard henselization fact recorded by the author: such etale extensions have the same henselization and the same completion as the pair (C, IC)
\end{itemize}\medskip
\ifdefined\hypertarget\hypertarget{sourcepointer-1}{}\fi
\noindent\textit{Source reference.} Original lemma 0A02: sources/stacks-chow/tags.txt, source lines 12620--12620. Prerequisite pointer invoked at 1465; it names a standard fact of a sibling chapter that the counted argument uses, and every step of the counted proof built on it is retained verbatim in the delivered body.\par\medskip
\ifdefined\hypertarget\hypertarget{sourcepointer-2}{}\fi
\noindent\textit{Source reference.} Original lemma 0AGV: sources/stacks-chow/tags.txt, source lines 13208--13208. Prerequisite pointer invoked at 1461; it names a standard fact of a sibling chapter that the counted argument uses, and every step of the counted proof built on it is retained verbatim in the delivered body.\par\medskip
\begin{lemma}
\label{lemma-algebraize-easy}
Let $I$ be an ideal of a Noetherian ring $A$. Let $C^h$ be the henselization
of a finite type $A$-algebra $C$ with respect to the ideal $IC$. Let
$J \subset C^h$ be an ideal. Then there exists a finite type $A$-algebra
$B$ such that $B^\wedge \cong (C^h/J)^\wedge$.
\end{lemma}

\begin{proof}
By More on Algebra, Lemma
\ref{more-algebra-lemma-henselization-Noetherian-pair}
the ring $C^h$ is Noetherian. Say $J = (g_1, \ldots, g_m)$.
The ring $C^h$ is a filtered colimit of \'etale $C$ algebras $C'$
such that $C/IC \to C'/IC'$ is an isomorphism
(see proof of More on Algebra, Lemma \ref{more-algebra-lemma-henselization}).
Pick an $C'$ such that $g_1, \ldots, g_m$ are the
images of $g'_1, \ldots, g'_m \in C'$.
Setting $B = C'/(g'_1, \ldots, g'_m)$ we get a finite
type $A$-algebra. Of course $(C, IC)$ and $C', IC')$ have
the same henselizations and the same completions.
It follows easily from this that $B^\wedge = (C^h/J)^\wedge$.
\end{proof}

\end{document}
